\section{Introduction}

Vulnerability prioritization is not a one-time decision. Enterprise operations teams schedule patching continuously across bi-weekly, monthly, or quarterly maintenance windows. Each window remediates a bounded set of (host, CVE) pairs, and the next window must prioritize from the residual backlog --- a backlog whose composition has changed because high-priority pairs have been addressed and new CVEs have been disclosed.

This temporal dimension is systematically absent from single-window evaluations, including the four prior papers in this research sequence. VulnPrio~\cite{paper1vulnprio} established a null result for CVE-level feature augmentation under a single synthetic window. HygienePrio~\cite{paper4hygieneprio} demonstrated that hygiene-augmented scoring achieves approximately 31 percentage-point Precision@50 improvement over EPSS-only at Window~1. Neither study addressed the question: \textit{does the advantage persist as the fleet evolves?}

The question matters operationally for two reasons. First, EPSS~\cite{jacobs2021epss} is designed as a CVE-level exploit-likelihood signal that is most informative when many high-EPSS CVEs remain unpatched. As high-priority CVEs are remediated, the remaining backlog increasingly consists of lower-EPSS items, and EPSS-only ranking degrades toward random ordering on the residual backlog. Second, hygiene signals --- patch posture, Active Directory (AD) exposure, telemetry freshness --- evolve on slower timescales than CVE remediation: a host with 60\% patch debt may still have 45\% debt three windows later if its patch velocity is low. This asymmetry suggests that hygiene-augmented scoring may be more temporally resilient than CVE-level scoring.

This paper tests this hypothesis through a pre-registered multi-window simulation extending HygienePrio's evaluation infrastructure to six consecutive bi-weekly windows.

\subsection{Contributions}

\begin{itemize}
\item \textbf{C1 --- Multi-window evaluation framework.} An extension of the EEHDA synthetic fleet simulation to support temporal fleet state evolution: patch application, new CVE disclosure, EPSS score update, and telemetry freshness accumulation across consecutive maintenance windows.

\item \textbf{C2 --- Temporal stability finding.} HygienePrio-full outperforms EPSS-only in all 150 of 150 window-seed evaluation pairs. EPSS-only decays from $0.331$ at Window~1 to $0.034$ at Window~6 (89.7\% relative drop); HygienePrio-full maintains P@50 $\geq 0.499$ at every window.

\item \textbf{C3 --- Operationally significant reversal.} By Window~2, EPSS-only's P@50 ($0.149$) has fallen below HRS-only ($0.229$), and HRS-only dominates EPSS-only at every subsequent window. Even the simplest hygiene baseline outperforms exploit-likelihood scoring once the second maintenance window begins.

\item \textbf{C4 --- Stability analysis.} HygienePrio-full shows lower P@50 variance across windows (more stable scheduling decisions) than EPSS-only, which exhibits a characteristic collapse pattern.
\end{itemize}

\subsection{Relationship to Prior Papers}

This paper is the fifth in the VulnPrio research sequence. VulnPrio~\cite{paper1vulnprio} reported a null result --- CVE-level augmentation $\approx$ EPSS-only at a single window. HygieneBench~\cite{paper3hygienebench} (HygieneBench) showed that hygiene signals carry discriminative structure in anomaly detection tasks. HygienePrio~\cite{paper4hygieneprio} (HygienePrio) demonstrated approximately 31 percentage-point Precision@50 improvement from hygiene augmentation at a single window. This paper extends HygienePrio to a temporal setting, showing that the hygiene advantage persists across six windows and that the absolute gap to EPSS-only grows from $26.4$ pp at Window~1 to $46.5$ pp at Window~6 as EPSS-only decays.
